\documentclass[../../main.tex]{subfiles}
\begin{document}

\begin{center}
  {\fontsize{48pt}{58pt}\selectfont \textbf{1974}\normalsize\textbf{年度}}
\end{center}

\begin{table}[htb]
  \centering
  \caption{出題テーマと難易度}
  \begin{tabular}{|c|p{9cm}|c|}
    \hline
    \textbf{問題番号} & \textbf{テーマ} & \textbf{難易度} \\
    \hline
    第1問 & 整数の累乗の一の位 & \\
    \hline
    第2問 & 放物線上の線分の中点の最小高さ & \\
    \hline
    第3問 & 投影図と空間内の直線の跡 & \\
    \hline
    第4問 & 直交する接線と導関数 & \\
    \hline
    第5問 & 円板の転がりによる点の軌跡 & \\
    \hline
    第6問 & 三勝先取の確率と試合数の期待値 & \\
    \hline
  \end{tabular}
\end{table}

\begin{table}[htb]
  \centering
  \caption{解答}
  \renewcommand{\arraystretch}{1.4}
  \begin{tabular}{|c|c|p{9cm}|}
    \hline
    \textbf{問題番号} & \textbf{小問} & \textbf{答え} \\
    \hline
    第1問 & (1) & $0,1,4,5,6,9$ \\
    \cline{2-3}
    & (2) & $f_5(n)=f_1(n)$（証明問題） \\
    \cline{2-3}
    & (3) & $0,1,5,6$ \\
    \hline
    第2問 & -- & $M=\left(\pm\dfrac{\sqrt{l-1}}{2},\dfrac{2l-1}{4}\right)$ \\
    \hline
    第3問 & -- & $C=\left(\sqrt{2}-2,\pm2^{3/4},0\right)$ \\
    \hline
    第4問 & -- & $a=1$ \\
    \hline
    第5問 & (1) & $(X,Y)=\left(4\cos^3\theta,6\sin\theta-4\sin^3\theta\right)$ \\
    \cline{2-3}
    & (2) & 最大値は$2\sqrt{2}$ \\
    \cline{2-3}
    & (3) & 弧の長さは$6$ \\
    \hline
    第6問 & (1) & $P-Q>p-q$ \\
    \cline{2-3}
    & (2) & $p=\dfrac{1}{2}\left(1+\sqrt{1-\dfrac{2\sqrt{30}}{15}}\right)$ \\
    \cline{2-3}
    & (3) & $p=\dfrac{1}{2}$で最大期待値$\dfrac{33}{8}$ \\
    \hline
  \end{tabular}
\end{table}

\paragraph{試験形態}
6問．試験時間は要確認．

\end{document}
